\fichatitulo{18 · CORPUS}{Texto para Discussão 357 · 2026}{Letras graúdas}
{\small Blanco. \textit{Letras gra{\'u}das, letras mi{\'u}das: leitura e complexidade textual dos discursos em Plen{\'a}rio 2007-2024}. Texto para Discussão 357 · 2026. \href{https://www12.senado.leg.br/publicacoes/estudos-legislativos/tipos-de-estudos/textos-para-discussao/td357}{Fonte consultada}.\par}

\campo{Problema e objetivo}
Como variaram características linguísticas dos discursos do Senado?

\campo{Principal contribuição · síntese interpretativa}
Produz uma descrição longitudinal de legibilidade e composição morfossintática, interpretando mudanças no registro discursivo.

\campo{Método / natureza da evidência}
Analisa 37.455 discursos com mais de 300 palavras, de 2007 a 2024; exclui apartes e aplica processamento linguístico com spaCy.

\campo{Resultado ou argumento principal}
Relata redução da complexidade medida e mudanças na distribuição de classes gramaticais; propõe a interpretação de um registro conversacional mediado.

\campo{Excertos-chave · seleção literal}
\begin{itemize}
\excerto{1}{mais de 300 palavras}{Procedimentos de seleção.}
\excerto{2}{conversacional mediado}{Conclusão.}
\end{itemize}

\campo{Limitações e análise crítica}
Análise da ficha: legibilidade calculada não mede diretamente compreensão cidadã. O limiar de extensão e a normalização condicionam o conjunto analisado; associações temporais não isolam causas.

\campo{Aproveitamento na dissertação · proposta}
Descrição dos dados: explicitar filtros, transcrições, extensão e possíveis diferenças entre períodos.

\registro{Fonte consultada nesta preparação: Resumo, procedimentos de seleção e conclusão. Chave: \texttt{blanco2026letrasGraudas}.}
